\documentclass[11pt,a4paper]{article}
\usepackage[margin=2.5cm]{geometry}
\usepackage{booktabs}
\usepackage{amssymb}
\usepackage{listings}
\usepackage{xcolor}
\usepackage{hyperref}
\usepackage{enumitem}
\usepackage{titlesec}

\definecolor{codebg}{RGB}{245,245,245}
\definecolor{codeframe}{RGB}{200,200,200}
\definecolor{statusdone}{RGB}{0,128,0}
\definecolor{statuspending}{RGB}{200,80,0}

\lstset{
  backgroundcolor=\color{codebg},
  frame=single,
  rulecolor=\color{codeframe},
  basicstyle=\ttfamily\small,
  breaklines=true,
  columns=fullflexible,
}

\newcommand{\itemdone}{\textcolor{statusdone}{$\checkmark$ Done}}
\newcommand{\itempending}{\textcolor{statuspending}{$\circ$ Pending}}

\title{\textbf{MoStream CloudLab Deployment Notes}}
\author{Nam Do}
\date{\today}

\begin{document}
\maketitle
\tableofcontents
\newpage

% ============================================================
\section{Cluster Layout}

\begin{center}
\begin{tabular}{llll}
  \toprule
  Node & Hostname & Internal IP & Role \\
  \midrule
  \texttt{apt051} & \texttt{node0.mostream-test.\ldots} & \texttt{10.10.1.1} & Kafka Broker \\
  \texttt{apt041} & \texttt{node1.mostream-test.\ldots} & \texttt{10.10.1.2} & MDWorkflow (PyFlink) \\
  \texttt{apt063} & \texttt{node2.mostream-test.\ldots} & \texttt{10.10.1.3} & TaskManager (spare) \\
  \texttt{apt055} & \texttt{node3.mostream-test.\ldots} & \texttt{10.10.1.4} & WLGenerator (Simulator) \\
  \bottomrule
\end{tabular}
\end{center}

SSH aliases (in \texttt{\textasciitilde/.ssh/config} on local machine):
\begin{lstlisting}
Host kafka
    HostName apt051.apt.emulab.net
    User NamSDSU
Host jobmanager
    HostName apt041.apt.emulab.net
    User NamSDSU
Host taskmanager
    HostName apt063.apt.emulab.net
    User NamSDSU
Host simulator
    HostName apt055.apt.emulab.net
    User NamSDSU
\end{lstlisting}

% ============================================================
\section{Node 0 --- apt051 (Kafka)}

\subsection{One-time fixes}
Fix hostname resolution:
\begin{lstlisting}
sudo sh -c 'echo "10.10.1.1 node0.mostream-test.pisceslabsd-pg0.apt.emulab.net" >> /etc/hosts'
\end{lstlisting}

\subsection{Start Kafka}
\begin{lstlisting}
cd ~/MoStream
KAFKA_HOST=10.10.1.1 docker compose up -d
\end{lstlisting}

Verify:
\begin{lstlisting}
docker compose ps
\end{lstlisting}

\subsection{Create topics (after every Kafka restart)}
\begin{lstlisting}
docker compose exec kafka kafka-topics --create \
  --topic Simulation --bootstrap-server localhost:9092 \
  --replication-factor 1 --partitions 1

docker compose exec kafka kafka-topics --create \
  --topic Result --bootstrap-server localhost:9092 \
  --replication-factor 1 --partitions 1

# Verify
docker compose exec kafka kafka-topics --list --bootstrap-server localhost:9092
\end{lstlisting}

\subsection{Monitor messages}
\begin{lstlisting}
docker compose exec kafka kafka-console-consumer \
  --bootstrap-server localhost:9092 \
  --topic Simulation --from-beginning --max-messages 5
\end{lstlisting}

% ============================================================
\section{Node 1 --- apt041 (MDWorkflow / PyFlink)}

\subsection{One-time fixes}

Fix hostname resolution (required for PyFlink MiniCluster):
\begin{lstlisting}
sudo sh -c 'echo "10.10.1.2 node1.mostream-test.pisceslabsd-pg0.apt.emulab.net" >> /etc/hosts'
\end{lstlisting}

Fix DNS (\texttt{systemd-resolved} not running on this node):
\begin{lstlisting}
sudo bash -c 'echo "nameserver 8.8.8.8" > /etc/resolv.conf'
\end{lstlisting}

\subsection{Package installs}
\begin{lstlisting}
source ~/miniconda3/etc/profile.d/conda.sh && conda activate mostream

# h5py 3.14 cannot read the model.h5 (non-standard float type)
pip install h5py==3.1.0

# nfp must have GlobalUpdate + ConcatDense (saved with Keras 2.8.0)
pip install nfp==0.1.3
\end{lstlisting}

\subsection{Data files required}
Copy model weights from Windows share:
\begin{lstlisting}
# On local machine:
ssh jobmanager "sudo mkdir -p /mnt/media/MDStream/StreamML/networks && \
  sudo chmod 777 /mnt/media/MDStream/StreamML/networks"

scp "/mnt/n/Applications/MoStreamData/.../model.h5" \
  jobmanager:/mnt/media/MDStream/StreamML/networks/
\end{lstlisting}

\subsection{Code fixes applied}
\textbf{MDWorkflow.py} --- Flink network memory (min/max must match):
\begin{lstlisting}
config.set_string("taskmanager.memory.network.min", "512m")
config.set_string("taskmanager.memory.network.max", "512m")  # was "2g"
\end{lstlisting}

\textbf{NPMMModel.py} --- add all required custom objects:
\begin{lstlisting}
custom_objects = nfp.custom_objects.copy()
custom_objects['ReduceAtoms'] = ReduceAtoms
if hasattr(nfp, 'GlobalUpdate'):  custom_objects['GlobalUpdate']  = nfp.GlobalUpdate
if hasattr(nfp, 'EdgeUpdate'):    custom_objects['EdgeUpdate']    = nfp.EdgeUpdate
if hasattr(nfp, 'NodeUpdate'):    custom_objects['NodeUpdate']    = nfp.NodeUpdate
if hasattr(nfp, 'ConcatDense'):   custom_objects['ConcatDense']   = nfp.ConcatDense
\end{lstlisting}

\subsection{Run MDWorkflow}
\begin{lstlisting}
source ~/miniconda3/etc/profile.d/conda.sh && conda activate mostream
cd ~/MoStream/MoStream/MDStream/StreamML

python MDWorkflow.py \
  --kafka-bootstrap 10.10.1.1:9092 \
  --starting-offset earliest \
  --group-id debug-group-1 2>&1 | tee output.log
\end{lstlisting}

% ============================================================
\section{Node 2 --- apt055 (WLGenerator / Simulator)}

\subsection{One-time fixes}

Create stub \texttt{gitinfo.py} (missing dependency):
\begin{lstlisting}
cat > ~/MoStream/MoStream/WLGenerator-node1/gitinfo.py << 'EOF'
def get_git_info():
    return {'commit': 'unknown'}
EOF
\end{lstlisting}

\subsection{Package installs}
\begin{lstlisting}
source ~/miniconda3/etc/profile.d/conda.sh && conda activate mostream
pip install requests qcelemental
\end{lstlisting}

\subsection{Data files required}
Copy training dataset from Windows share:
\begin{lstlisting}
# On local machine:
ssh simulator "sudo mkdir -p /mnt/media/MDStream/WLGenerator/dataset && \
  sudo chmod 777 /mnt/media/MDStream/WLGenerator/dataset"

scp "/mnt/n/Applications/MoStreamData/.../training-data-simple.txt" \
  simulator:/mnt/media/MDStream/WLGenerator/dataset/
\end{lstlisting}

\subsection{Run simulator}
\begin{lstlisting}
source ~/miniconda3/etc/profile.d/conda.sh && conda activate mostream
cd ~/MoStream/MoStream/WLGenerator-node1

KAFKA_BOOTSTRAP=10.10.1.1:9092 python simulator.py \
  --topic Simulation --interval 1.0
\end{lstlisting}

% ============================================================
\section{Node 3 --- apt063 (TaskManager / Spare)}
Not yet configured. No steps needed for current single-node PyFlink setup.

% ============================================================
\section{Progress Summary}

\begin{tabular}{p{6cm}ll}
  \toprule
  Item & Status & Notes \\
  \midrule
  Kafka running on apt051           & \itemdone    & Topics wiped on restart \\
  Advertised listener fix           & \itemdone    & \texttt{KAFKA\_HOST=10.10.1.1} \\
  Simulator sending messages        & \itemdone    & Set \texttt{KAFKA\_BOOTSTRAP} env var \\
  MDWorkflow receiving + training   & \itemdone    & \texttt{train\_y} accumulates correctly \\
  Flink network memory fix          & \itemdone    & min = max = 512m \\
  h5py downgrade to 3.1.0           & \itemdone    & Newer versions break model.h5 load \\
  nfp pinned to 0.1.3               & \itemdone    & Needs GlobalUpdate + ConcatDense \\
  NPMMModel custom\_objects fix     & \itempending & ConcatDense added, not yet tested \\
  Full inference + ranking          & \itempending & Blocked on above \\
  Results in Kafka Result topic     & \itempending & Blocked on inference \\
  apt063 TaskManager setup          & \itempending & Not started \\
  \bottomrule
\end{tabular}

\bigskip
\textbf{Known issues:}
\begin{itemize}
  \item Kafka topics are wiped on \texttt{docker compose down} --- recreate them every restart.
  \item \texttt{systemd-resolved} is broken on apt041 --- \texttt{/etc/resolv.conf} fix is not persistent across reboots; re-apply if needed.
  \item Model cold-start: NPMM returns \texttt{model\_not\_ready} until 16 messages are accumulated.
  \item No Flink checkpointing --- trained weights lost if MDWorkflow crashes.
\end{itemize}

\end{document}